% from ega2-1-fr.tex
\oldpage[II]{135}

$g^{-1}(X_i)$ de $X'$, $g_i:X'_i\to X_i$ le morphisme $g$ restreint à
$X'_i$, qui est fini (\hyperref[II.6.1.5-fr]{6.1.5, (iii)}) et surjectif~;
soit $k_i$ le \emph{rang} de la $\mathcal{O}_{X_i}$-Algèbre
$\mathcal{B}_i=(g_i)_*(\mathcal{O}_{X'_i})$. Comme
$\mathcal{B}_i\otimes_{\mathcal{O}_{X_i}}\mathcal{R}(X_i)$ est un
préfaisceau constant (\hyperref[I.7.3.5-fr]{I, 7.3.5}), le rang $k_i$ est
aussi le rang de la $(\mathcal{O}_X|U)$-Algèbre $\mathcal{B}|U$ pour tout
ouvert $U$ de $X$ ne rencontrant que la \emph{seule} composante irréductible
$X_i$. Si $T$ est un ensemble ouvert dans $X$ tel que $\mathcal{B}|T$ soit
isomorphe à $\mathcal{O}_X^m|T$, il résulte de la remarque précédente que les
nombres $k_i$ sont égaux à $m$ pour tous les indices $i$ tels que
$T\cap X_i\neq\varnothing$. Soit alors $U$ l'ouvert de $X$ formé des points
au voisinage desquels $\mathcal{B}$ est un $\mathcal{O}_X$-Module localement
libre, et soient $U_j$ ($1\leq j\leq s$) ses composantes connexes, qui sont
ouvertes dans $X$ et en nombre fini (puisque $U$ est noethérien)~; désignons
par $V_j$ le sous-préschéma fermé de $X'$, \emph{adhérence} du sous-préschéma
induit sur l'ouvert $g^{-1}(U_j)$
(\hyperref[I.9.5.11-fr]{I, 9.5.11}). D'après ce qui précède, pour tous les
indices $i$ tels que $X_i\cap U_j\neq\varnothing$, les rangs $k_i$ sont tous
égaux à un même entier $m_j$~; on notera d'ailleurs qu'un même $X_i$ ne peut
rencontrer deux $U_j$ d'indices distincts. Soient $i_\lambda$ les indices
$i$ tels que $X_i\cap U=\varnothing$. Considérons le produit $k$ de tous les
$k_i$, posons $n_i=k/k_i$, et soit $X''$ le préschéma défini comme suit.
Pour chaque $j$ ($1\leq j\leq s$), on considère $k/m_j$ préschémas
isomorphes à $V_j$, et pour chaque $\lambda$, $k/k_{i_\lambda}$ préschémas
isomorphes à $X'_{i_\lambda}$~; $X''$ est la \emph{somme} de tous ces
préschémas. On définit un morphisme $g'':X''\to X'$, se réduisant à
l'injection canonique dans chacun des summandes de $X''$~; il est clair que
$g''$ est un morphisme fini dominant, donc surjectif (puisqu'un morphisme
fini est fermé (\hyperref[II.6.1.10-fr]{6.1.10}))~; posons
$g'=g\circ g''$, qui est un morphisme fini surjectif $X''\to X$~; on a
${g'}^*(\mathcal{L})={g''}^*(g^*(\mathcal{L}))$, donc
${g'}^*(\mathcal{L})$ est un $\mathcal{O}_{X''}$-Module ample
(\hyperref[II.5.1.12-fr]{5.1.12}). Il est clair alors que pour ce nouveau
préschéma $X''$, les rangs définis comme les $k_i$ pour $X'$ sont tous
\emph{égaux à $k$}~; compte tenu de
(\hyperref[I.7.3.3-fr]{I, 7.3.3}), on en conclut aussitôt que pour tout ouvert
affine $T$ de $X$,
\[
\bigl(g'_*(\mathcal{O}_{X''})
\otimes_{\mathcal{O}_X}\mathcal{R}(X)\bigr)|T
\]
est un $(\mathcal{R}(X)|T)$-Module isomorphe à
$(\mathcal{R}(X)|T)^k$.

\begin{corollary}[6.6.12]
\label{II.6.6.12-fr}
Si, dans l'énoncé de (\hyperref[II.6.6.11-fr]{6.6.11}), on a $W=X$, alors
pour qu'un $\mathcal{O}_X$-Module inversible $\mathcal{L}$ soit ample
relativement à $f$, il faut et il suffit que $g^*(\mathcal{L})$ soit ample
relativement à $f\circ g$.
\end{corollary}

\begin{remark}[6.6.13]
\label{II.6.6.13-fr}
Au chapitre III, nous verrons que si $Y$ est noethérien, $f$ de type fini, et
si la restriction de $f$ au sous-préschéma fermé réduit de $X$ ayant $X-W$
pour espace sous-jacent est \emph{propre}, alors la conclusion de
(\hyperref[II.6.6.12-fr]{6.6.12}) est encore valable. Mais nous donnerons au
chapitre V des exemples de schémas algébriques $X$ sur un corps $K$ (le
morphisme structural $X\to\operatorname{Spec}(K)$ n'étant pas propre) dont
le normalisé $X'$ est quasi-affine, mais qui n'est pas quasi-affine (de sorte
que $\mathcal{O}_X$ n'est pas ample, bien que $\mathcal{O}_{X'}$ le soit
(\hyperref[II.5.1.12-fr]{5.1.12}) et que le morphisme $X'\to X$ soit fini
surjectif (\hyperref[II.6.3.10-fr]{6.3.10})). Nous allons voir au numéro
suivant que cette circonstance ne peut se produire lorsqu'on remplace
«~quasi-affine~» par «~affine~».
\end{remark}

\subsection{Le théorème de Chevalley.}
\label{subsection:II.6.7-fr}

Nous allons établir (à l'aide du critère de Serre
(\hyperref[II.5.2.1-fr]{5.2.1})) le théorème suivant, qui a été démontré par
C.~Chevalley par d'autres méthodes, dans le cas des schémas
\emph{algébriques}.

